% BHU PYQ -- Transcribed & Corrected
% Paper: CHEMN41 : Physical Chemistry-II (End-Term)
% Exam: B.Sc. (Hons) Semester IV Examination, 2025-26 (NEP)
% Department: Chemistry

\documentclass[11pt,a4paper]{article}
\usepackage[a4paper,top=2.5cm,bottom=2.5cm,left=2.8cm,right=2.8cm,headheight=14pt]{geometry}
\usepackage{amsmath,amssymb}
\usepackage[T1]{fontenc}
\usepackage[utf8]{inputenc}
\usepackage{lmodern,microtype,fancyhdr,enumitem,hyperref,lastpage,parskip}

\hypersetup{
    colorlinks=true,
    urlcolor=black,
    linkcolor=black,
    pdftitle={CHEMN41: Physical Chemistry-II (End-Term) -- BHU Sem IV 2025-26 (NEP)}
}

\pagestyle{fancy}
\fancyhf{}
\renewcommand{\headrulewidth}{0.4pt}
\renewcommand{\footrulewidth}{0.4pt}
\fancyhead[L]{\small   \textit{Banaras Hindu University}}
\fancyhead[R]{\small   \textit{CHEMN41: Physical Chemistry-II (2025-26)}}
\fancyfoot[C]{\small Page  \thepage\ of \pageref{LastPage}}


\newlist{parts}{enumerate}{2}
\setlist[parts,1]{label=(\alph*),leftmargin=*,itemsep=4pt,topsep=2pt}
\setlist[parts,2]{label=(\roman*),leftmargin=*,itemsep=2pt,topsep=2pt}

\newcommand{\pts}[1]{\hfill   \textit{[#1]}}


\begin{document}


\begin{center}
  {\large\bfseries BANARAS HINDU UNIVERSITY}\\[2pt]
  {
\normalsize B.Sc. (Hons.) --- Semester IV (NEP) Examination, 2025-26}\\[6pt]
  \rule{0.8\linewidth}{0.5pt}\\[6pt]
  {\large\bfseries CHEMISTRY}\\[4pt]
  {\large\bfseries Paper Code: CHEMN41 : Physical Chemistry-II}\\[4pt]
  {
\normalsize Time: 2.5 Hours\hfill Maximum Marks: 70}\\[2pt]
  {\small REF NO. Conf/Even/N/2025-26/08}
\end{center}

\rule{\linewidth}{0.4pt}

\smallskip



\noindent   \textit{Instructions: Answer   \textbf{five} questions including Question No.~1 which is compulsory. All questions carry equal marks (14 marks each).}

\bigskip




\noindent   \textbf{Question 1.}   \textit{(Compulsory, 14 Marks Total)}

\begin{parts}
  \item Define the quantum yield of a photochemical reaction.\pts{2}
  \item What do you mean by fugacity?\pts{2}
  \item Write the rate equation of a unimolecular surface reaction, identify the terms involved and provide one example.\pts{2}
  \item Write the law of photochemical equivalence and as per this law, what should be the quantum yield?\pts{2}
  \item What are the general characteristics of radioactive decay?\pts{2}
  \item How will you synthesise radioactive sodium in the laboratory?\pts{2}
  \item What is artificial radioactivity? Give an example.\pts{2}
\end{parts}

\medskip\rule{\linewidth}{0.2pt}\medskip




\noindent   \textbf{Question 2.}

\begin{parts}
  \item Write the expression to find out the Gibbs energy of mixing ($\Delta G_{ \text{mix}}$) and the entropy of mixing ($\Delta S_{ \text{mix}}$).\pts{3}
  \item Write van't Hoff equation for Osmotic pressure and discuss how the molecular weight of a polymer can be measured through osmotic pressure measurement?\pts{4}
  \item Derive Gibbs-Duhem equation and explain the significance of the Gibbs-Duhem equation in multi-component systems.\pts{7}
\end{parts}

\medskip\rule{\linewidth}{0.2pt}\medskip




\noindent   \textbf{Question 3.}

\begin{parts}
  \item Write any two deactivation processes of a photochemically activated molecule.\pts{3}
  \item What do you mean by primary and secondary processes of photochemical reaction? And mention the process responsible for the variation of the quantum yield.\pts{4}
  \item Write the Brunauer-Emmett-Teller (BET) adsorption isotherm equation. Using the linear form of the BET equation, explain how it can be used to determine the specific surface area of the solid.\pts{7}
\end{parts}

\medskip\rule{\linewidth}{0.2pt}\medskip




\noindent   \textbf{Question 4.}

\begin{parts}
  \item What are the two differences between thermal equilibrium and photostationary state?\pts{3}
  \item What is Michaelis-Menten equation? Draw the Lineweaver-Burk plot (double reciprocal plot) and write the values of the slope and intercept of the plot.\pts{4}
  \item Discuss the fluorescence quenching, considering the dimerization of anthracene and derive the Stern-Volmer equation.\pts{7}
\end{parts}

\medskip\rule{\linewidth}{0.2pt}\medskip




\noindent   \textbf{Question 5.}

\begin{parts}
  \item Describe the Duhem-Margules equation with one of its applications.\pts{3}
  \item Depicting the graph, discuss the influence of the pH and the temperature on the enzyme catalysed reaction.\pts{4}
  \item Determine the rate of photochemical decomposition of $ \text{HI}$ (showing the reaction mechanism) and the quantum yield of this process.\pts{7}
\end{parts}

\medskip\rule{\linewidth}{0.2pt}\medskip




\noindent   \textbf{Question 6.}

\begin{parts}
  \item Describe the concept of Bohr's compound nucleus theory of nuclear reactions.\pts{3}
  \item What are radioisotopes? Write the precautions to be taken in the use of radioisotopes as tracer.\pts{4}
  \item Derive an expression for the decay constant of a radioactive substance. How is this decay constant related to half-life?\pts{7}
\end{parts}

\medskip\rule{\linewidth}{0.2pt}\medskip




\noindent   \textbf{Question 7.}

\begin{parts}
  \item What are the different types of nuclear reactions? Describe radiative capture and transfer reactions with suitable examples.\pts{6}
  \item Explain the direct isotopic dilution analysis and inverse isotopic dilution analysis methods.\pts{8}
\end{parts}

\medskip\rule{\linewidth}{0.2pt}\medskip




\noindent   \textbf{Question 8.} Write short notes on any   \textbf{two} of the following:\pts{$2  \times 7 = 14$}

\begin{parts}
  \item Osmosis and osmotic pressure
  \item Compton scattering
  \item Photochemical reactions between $ \text{H}_2$ and $ \text{Br}_2$ and derive thereof the rate of formation of $ \text{HBr}$.
\end{parts}

\vfill
\rule{\linewidth}{0.4pt}

\begin{center}\small
\href{https://bhu-exam.vercel.app/}{Click here to visit our website}\\[4pt]
\href{https://me-aryan.vercel.app/}{Click here to visit our contributor}
\end{center}

\end{document}
